%% secoes/introducao/02-pre-termo.tex — incluído por main.tex

%\subsection{Nascimento pré-termo}\label{subsec:intro-pt}

Já o nascimento pré-termo é definido quando ocorre antes de 37 semanas completas de gestação. A prematuridade é uma das principais causas de morte de crianças menores de cinco anos e está associada a efeitos de longo prazo sobre a saúde física, o neurodesenvolvimento e as condições socioeconômicas dos que sobrevivem \cite{ohuma2023}. Estima-se que 13,4 milhões de crianças nasceram pré-termo no mundo em 2020, 9,9\% de todos os nascimentos, sem mudança mensurável em relação a 2010 \cite{ohuma2023}.

%No Brasil, a prevalência de nascimentos pré-termo foi de 11,1\% entre 2011 e 2021 \cite{alberton2023} e passou de 11,3\% em 2014 para 11,9\% em 2023 \cite{victor2025}, acima da média mundial. As maiores prevalências recaem sobre mulheres indígenas e pretas, com menor escolaridade, nas idades extremas e com poucas consultas de pré-natal, e sobre as regiões Norte e Nordeste \cite{alberton2023,victor2025}.

%Na população de baixa renda, o número reduzido de consultas de pré-natal está entre os fatores de maior risco, tanto para o primeiro nascimento pré-termo quanto para a sua recorrência \cite{rocha2022}. As políticas de proteção social também se relacionam ao desfecho: entre mulheres de baixa renda, o recebimento do Bolsa Família durante a gestação esteve associado à redução dos nascimentos pré-termo extremos \cite{ortelan2024}.
